\documentclass{article}
\usepackage[T1]{fontenc}
\usepackage{lmodern}
\usepackage{graphicx}
\usepackage{amsmath,amssymb}
\usepackage[a4paper,margin=2cm]{geometry}
\usepackage[absolute,overlay]{textpos}
\setlength{\TPHorizModule}{1mm}
\setlength{\TPVertModule}{1mm}
\usepackage[indonesian]{babel}
\usepackage{hyperref}
\setlength{\emergencystretch}{2em}
% unit: Halaman 35

\begin{document}

% === Halaman 35 (python/native) ===

35

• Contoh: jika LFSR 4-bit diinisialisasi dengan U = 1111

• Barisan bit acak: 1 1 1 1 0 1 0 1 1 0 0 1 0 0 0 …

• Periode LFSR n-bit: 2n – 1 

i 


Isi Register 


Bit Keluaran 





0 


1 1 1 1 





1 


0 1 1 1 



1 


2 


1 0 1 1 



1 


3 


0 1 0 1 



1 


4 


1 0 1 0 



1 


5 


1 1 0 1 



0 


6 


0 1 1 0 



1 


7 


0 0 1 1 



0 


8 


1 0 0 1 



1 


9 


0 1 0 0 



1 


10 


0 0 1 0 



0 


11 


0 0 0 1 



0 


12 


1 0 0 0 



1 


13 


1 1 0 0 



0 


14 


1 1 1 0 



0 

15                            1   1   1   1                                               0

\end{document}
